\newcommand{\DoNotLoadEpstopdf}{}
\documentclass[11pt,reqno]{amsart}
\usepackage[T1]{fontenc}
\usepackage{lmodern,amsmath,amssymb,amsthm,mathtools,mathrsfs,microtype}
\usepackage[margin=1.1in]{geometry}
\usepackage[hidelinks]{hyperref}
\hypersetup{pdftitle={Collision records, arithmetic, and raw local inversion in a triangular Lorentz gas},pdfauthor={Qian Qi}}
\setlength{\emergencystretch}{2em}
\allowdisplaybreaks[2]
\numberwithin{equation}{section}
\newtheorem{theorem}{Theorem}[section]
\newtheorem{lemma}[theorem]{Lemma}
\newtheorem{proposition}[theorem]{Proposition}
\newtheorem{corollary}[theorem]{Corollary}
\theoremstyle{definition}
\newtheorem{definition}[theorem]{Definition}
\newtheorem{inputstatement}[theorem]{Input statement}
\theoremstyle{remark}\newtheorem{remark}[theorem]{Remark}
\newcommand{\R}{\mathbb R}\newcommand{\Z}{\mathbb Z}
\newcommand{\one}{\mathbf1}\newcommand{\dd}{\,\mathrm d}
\newcommand{\esssup}{\operatorname*{ess\,sup}}
\newcommand{\diver}{\operatorname{div}}
\newcommand{\TV}{\operatorname{TV}}
\title[Collision records and raw local inversion]{Collision records, arithmetic, and raw local inversion in a triangular Lorentz gas}
\author{Qian Qi}
\date{October 11, 2026. A2-DYN, revision 77}
\subjclass[2020]{37D50, 37A50, 60F05, 28A75}
\keywords{Sinai billiards, exact collision records, local limits, endpoint momenta, coarea transport, optical Jacobi matrices}
\begin{document}
\raggedbottom
\begin{abstract}
We retain the pointwise local limit problem for the full displacement,
collision count, occupation, and roof record of actual returns in a
triangular Lorentz gas. The exact initial and terminal section tests
give collision-uniform positive bounds for the Hessian in physical
endpoint angles. They determine an explicit roof flow preserving the
angle of the endpoint momenta. Its forward motion is uniformly
square-root controlled and never reaches an interior critical point;
no Morse collars or cap-dependent patching are needed. We disintegrate
the complete original source jointly over roof and momentum angle,
and construct a maximal angle-preserving subcoupling after summing
all actual words with the same exact label. This recoupling can recover
mass lost by an individual stopped orbit. A finite signed angular
current retains true boundary flux and the positive critical residues.
These results sharpen the complete-source persistence estimate for
the original pointwise law. The collision-normalized concentration
bound for the complete angular defect remains the quantitative
estimate required to close that law.
\end{abstract}
\maketitle

\section{The unchanged physical problem}
Put $R_-=9/20$ and $R_+=47/100$, and let $R\in I=[R_-,R_+]$. Circular scatterers have centers in
$\Lambda=\Z(1,0)+\Z(1/2,\sqrt3/2)$. Write $T_R$ for the regular
collision map and $(\alpha,p)\in(\R/2\pi\Z)\times(-1,1)$ for its
coordinates; $p$ is the outgoing tangential velocity and
$\dd\nu=\dd\alpha\,\dd p/(4\pi)$. Physical flights satisfy
$\tau_-\le\tau_R\le\tau_+$, where $\tau_-=3/50$ and $\tau_+$
is the fixed finite-horizon bound.

The section is specified here to fix every event in the argument.
Put $h_0=1/200$. Let $Y_R$ be the union of the four open coordinate
squares of half-side $h_0$ centered at
\[
 (0,0),\quad(\pi/6,-1/2),\quad
 (-\theta_R,\sin\theta_R),\quad
 (\pi/3-\theta_R,\sin\theta_R).
\]
Here $\theta_R=2\arctan q_R$, and $q_R$ is the unique root in
$(49/100,53/100)$ of
\[
 8Rq(1+q^2)-4q(1-q^2)+\sqrt3(1-6q^2+q^4)=0.
\]
Set
\[
 U=\{|\alpha-\pi/6|<3/50,\ |p|<3/20\},\qquad
 Y_R^*=Y_R\cup U,\qquad c_*=\frac{91}{10000\pi}.
\]
The rectangles are disjoint and $\nu(Y_R^*)=c_*$. We always use
$\nu_R^*=c_*^{-1}\nu|_{Y_R^*}$, not a renormalization of a smaller
source.

At $m$ flights, let $L=L_{m,R}$ be their total length,
$K^{\rm c}_{m,R}$ their lattice displacement, and
$A_{m,R}=\sum_{j=0}^{m-1}\one_{Y_R^*}\circ T_R^j$. The event is
\begin{equation}\label{eq:v76-event}
 \Pi_{n,k,m,R}=Y_R^*\cap\{K^{\rm c}_{m,R}=k,\ A_{m,R}=n\}
                         \cap T_R^{-m}Y_R^*.
\end{equation}
The initial visit is counted; the terminal visit is tested separately.
On this event the actual $n$th return has collision count $m$ and
roof $L$. Denote its density by $p_{n,R}(k,m,t)$.

Write $c_j=\sqrt{1-p_j^2}$ and let $\Delta_j$ be the least clearance
of flight $j$ from any third disk, excluding its two incident disks.
Keep the original smooth function $0\le\Theta\le1$, zero on
$(-\infty,1]$ and one on $[2,\infty)$. With $q_0=47/53$,
$d_j=\min(j,m-j)$ and $d'_j=\min(j,m-1-j)$, set
\begin{align}
 H^\varepsilon&=
 \prod_{j=0}^{m}\Theta\!\left(\frac{c_j}{\varepsilon q_0^{d_j/4}}\right)
 \prod_{j=0}^{m-1}\Theta\!\left(\frac{\Delta_j}{\varepsilon q_0^{d'_j/4}}\right),
 \qquad W^\varepsilon=1-H^\varepsilon,\label{eq:v76-guard}\\
 b^\varepsilon(t)\dd t
   &=L_\#(W^\varepsilon\one_{\Pi_{n,k,m,R}}\nu_R^*).
 \label{eq:v76-source}
\end{align}
In particular first incidence and first clearance outside any selected
critical neighborhoods are still present. No section-distance factor
has been inserted into this guard.

\subsection{The arithmetic reference and the target}
For $y=(k_1,k_2,t,n)$ put
\[
 \mu_R=(0,0,\bar\tau_R,c_*),\quad
 Z=(y-m\mu_R)/\sqrt m,\quad
 \mathcal C_{m,R,n,k}(M)=\{t:|Z|\le M\}.
\]
We retain the finite spectral-peak reference, including peaks close to
but not on the unit circle. To specify its meaning, the collision
record observable is $(\kappa_1,\kappa_2,\tau,\one_{Y_R^*})$.
Let $\lambda_j(R,\vartheta)$ and $\mathsf P_j(R,\vartheta)$ be its
continued simple twisted eigenvalue and projection in each of the
fixed peak charts, and let $\chi_j$ be their frequency partition.
At the real maximum $a_j$ of $\log|\lambda_j|$ put
\[
 z_j=\lambda_j(R,a_j),\quad
 \mu_j=i^{-1}\nabla\log\lambda_j(R,a_j),\quad
 H_j=-\nabla^2\log\lambda_j(R,a_j),
\]
with $\operatorname{Re}H_j\ge aI_4$ for a common $a>0$. If
$\eta_R=\one_{Y_R^*}$ and $\ell$ denotes mass evaluation, put
$B_j=\chi_j(R,a_j)\ell(M_{\eta_R}\mathsf P_j(R,a_j)(\eta_R\nu))$.
Define
\begin{align}
 G_H(z)&=(2\pi)^{-4}\int_{\R^4}e^{-iv\cdot z-v^{\mathsf T}Hv/2}\dd v,
 \notag\\
 \mathcal L_{m,R}(y)&=\operatorname{Re}\left[
 c_*^{-1}\sum_j e^{-ia_j\cdot y}z_j^m B_j
      G_{H_j}\left((y-m\mu_j)/\sqrt m\right)\right].
 \label{eq:v76-reference}
\end{align}
Only the charts active at $R$ are summed. This is the same transition
reference as in the retained spectral theorem. The existence and
uniform estimates for these scalar data belong to that analytic
input, not to the transport argument below. No frequency branch or
arithmetic zero class is deleted.

The governing assertion remains
\begin{equation}\label{eq:v76-target}
 \sup_{R,\,1\le n\le m,\,k\in\Z^2}
 \esssup_{t\in\mathcal C_{m,R,n,k}(M)}
 |m^2p_{n,R}(k,m,t)-\mathcal L_{m,R}(k_1,k_2,t,n)|\longrightarrow0.
\end{equation}
Let $\mathcal E_M$ denote this error's collision limsup. An empty
central target contributes zero. The order of the uniformity is
$\sup_R\esssup_t$, so exceptional density sets may depend on $R$.

\subsection{The analytic inputs in the form used here}
\begin{inputstatement}\label{input:v76-analytic}
The unchanged source \eqref{eq:v76-source} has a nonnegative,
compactly supported $BV$ density at every fixed count. For every
fixed $\varepsilon,\delta>0$, with
$A_\delta f(t)=(2\delta)^{-1}\int_{t-\delta}^{t+\delta}f(s)\dd s$,
\begin{equation}\label{eq:v76-average}
 \limsup_m\sup_{R,n,k}m^2\|A_\delta b^\varepsilon\|_\infty
                                      \le C\varepsilon^{1/16}.
\end{equation}
The constant is independent of the previously fixed width. For large
fixed $B$ and $\varepsilon(B)=A_0B^{-1/12}$,
\begin{equation}\label{eq:v76-remainder}
 m^2p-\mathcal L=m^2b^{\varepsilon(B)}+e_B,\qquad
 \limsup_m\sup_{R,n,k}\esssup_{\mathcal C(M)}|e_B|
                                     \le C_MB^{-1/192}.
\end{equation}
All auxiliary spectral bands in these statements are fixed before
the collision limit and enlarged only afterwards.
\end{inputstatement}
These are the complete-source tame theorem, the fixed-width
complete-source bound, and the positive scalar remainder theorem in
the technical supplement. Their proofs and historical hypotheses are
retained there. The fixed-count source uses threshold integration
\cite{CMInt}, definable monotonicity \cite{Coste}, and the
distributional characterization of $BV$ \cite{AFP}; it does not
assume definability of the original smooth guard. They are restated as inputs to distinguish them from
the new geometric and transport proofs. No long-time bound on a
fixed-count variation constant is part of this statement.


\subsection{The new physical construction}
The retained convex-support and optical arguments provide genuine
information on every regular physical word. They do not on their
own control the complete weighted sum on an exact roof. The present
construction uses a different observable from the auxiliary cap
normal: the angle of the two physical endpoint momenta. The fixed
section makes its coarea Jacobian uniformly positive after the
optical twist is factored out. The resulting field is explicit and
requires no collision-dependent critical neighborhoods.

The new subcoupling is constructed after summing all actual words of
one exact event. It may match mass between different true words,
while preserving their exact displacement, occupation, terminal test,
and endpoint-momentum angle. This is more permissive than continuation
of a single physical orbit and more restrictive than unconstrained
scalar matching. Its finite-current estimate is an unconditional
fixed-count statement. Its collision-normalized concentration remains
a distinct dynamical question.

\input{core/174_canonical_endpoint_momentum_flow}
\input{core/175_complete_angular_recoupling}
\input{core/176_angular_current_and_ordered_endpoint}

\section{The remaining complete-source estimate}
The principal assertion remains \eqref{eq:v76-target}, with its full
arithmetic transition and all zero classes. The new uniform bounds
concern the actual endpoint Hessian and the forward motion of an
explicit physical field. They remove arbitrary field choice and
Morse-radius patching from that construction. The exact angular
sum and its maximal recoupling also avoid charging a separate loss
when another actual word supplies matching mass at the same angle
and label. True boundaries are neither declared null in flux nor
removed from the original guard.

What is still needed is \eqref{eq:v77-remaining}, or a direct estimate
of comparable strength for the original complete source. Neither
$\int A^\varepsilon\le1$ nor the finite constants $N_m,J_m,C_m$
control concentration at one roof after multiplication by $m^2$.
The explicit drift formula also involves third derivatives and source
derivatives for which no collision-uniform summed bound has been
obtained here. Thus the organizing pointwise theorem and its
unconditional path and likelihood consequences are not claimed as
proved. All earlier arguments, including their exact hypotheses,
are retained in the technical supplement; the target is not replaced
by a smoothed law or a different record.

\input{journal-references}
\end{document}
